% GUIA DE TRIAGEM — Volume 5 (gerado por expandir_volume.py)

\chapter{Guia de Observação e Encaminhamento}

**Este guia é de uso pedagógico.** No 5º ano, o professor observa os domínios abaixo ao longo das unidades e registra no Registro de Progresso. Sinais persistentes (2+ checkpoints com R/N no mesmo domínio) são discutidos com a família e avaliados pelo profissional habilitado com o Volume Profissional de Sondagem e Rastreio.

\section{Domínio: Leitura e Fluência}
Observar precisão, ritmo e prosódia na leitura oral de frases e textos curtos do volume.

\begin{center}
\renewcommand{\arraystretch}{1.4}
\begin{tabular}{@{}p{9cm}cccc@{}}
\toprule
Indicador observável & S & F & R & N \\
\midrule

Lê frases com precisão & & & & \\

Lê com ritmo adequado & & & & \\

Usa entonação ao ler & & & & \\

\bottomrule\end{tabular}\end{center}

\section{Domínio: Escrita e Ortografia}
Observar correspondência grafema-fonema, ortografia das grafias-alvo do ano e legibilidade do traçado.

\begin{center}
\renewcommand{\arraystretch}{1.4}
\begin{tabular}{@{}p{9cm}cccc@{}}
\toprule
Indicador observável & S & F & R & N \\
\midrule

Escreve com ortografia estável & & & & \\

Usa as grafias-alvo do ano & & & & \\

Traçado legível e homogêneo & & & & \\

\bottomrule\end{tabular}\end{center}

\section{Domínio: Consciência Fonológica}
Observar segmentação, aliteração, rima e manipulação de fonemas em palavras com as grafias do ano.

\begin{center}
\renewcommand{\arraystretch}{1.4}
\begin{tabular}{@{}p{9cm}cccc@{}}
\toprule
Indicador observável & S & F & R & N \\
\midrule

Segmenta palavras em sílabas & & & & \\

Percebe rimas e aliterações & & & & \\

Manipula fonemas (troca/omite) & & & & \\

\bottomrule\end{tabular}\end{center}

\section{Domínio: Linguagem Oral}
Observar vocabulário, narrativa oral, compreensão e articulação na fala espontânea.

\begin{center}
\renewcommand{\arraystretch}{1.4}
\begin{tabular}{@{}p{9cm}cccc@{}}
\toprule
Indicador observável & S & F & R & N \\
\midrule

Narra fatos com sequência lógica & & & & \\

Usa vocabulário variado & & & & \\

Articula sem trocas persistentes & & & & \\

\bottomrule\end{tabular}\end{center}

\section{Domínio: Atenção e Funções Executivas}
Observar manutenção de foco, planejamento da tarefa, inibição e memória de trabalho nas folhas.

\begin{center}
\renewcommand{\arraystretch}{1.4}
\begin{tabular}{@{}p{9cm}cccc@{}}
\toprule
Indicador observável & S & F & R & N \\
\midrule

Mantém foco por 10–15 minutos & & & & \\

Planeja a tarefa antes de executar & & & & \\

Inibe respostas impulsivas & & & & \\

\bottomrule\end{tabular}\end{center}

\section{Domínio: Motricidade Fina e Traçado}
Observar preensão do lápis, pressão, direção e planejamento motor na escrita cursiva.

\begin{center}
\renewcommand{\arraystretch}{1.4}
\begin{tabular}{@{}p{9cm}cccc@{}}
\toprule
Indicador observável & S & F & R & N \\
\midrule

Preensão funcional do lápis & & & & \\

Pressão adequada no papel & & & & \\

Direção correta do traçado & & & & \\

\bottomrule\end{tabular}\end{center}

\section{Domínio: Comportamento e Interação}
Observar regulação emocional, cooperação e resposta à frustração nas atividades coletivas.

\begin{center}
\renewcommand{\arraystretch}{1.4}
\begin{tabular}{@{}p{9cm}cccc@{}}
\toprule
Indicador observável & S & F & R & N \\
\midrule

Regula emoções com apoio & & & & \\

Coopera em duplas & & & & \\

Lida com a frustração & & & & \\

\bottomrule\end{tabular}\end{center}

\section{Fluxo de Encaminhamento}

1) Registrar no Registro de Progresso. 2) Conversar com a família (nunca dar diagnóstico). 3) Encaminhar ao profissional habilitado para aplicação das fichas do Volume Profissional. 4) Acompanhar e manter as rotas diferenciadas em sala.

\section{Modelo de Registro Individual do Aluno}

O registro individual acompanha cada aluno ao longo do ano. Preencha quinzenalmente, anexando as folhas Kumon e os checkpoints.

\begin{center}
\renewcommand{\arraystretch}{1.5}
\begin{tabular}{@{}p{3cm}p{3cm}p{5cm}p{3cm}@{}}
\toprule
Período & Foco observado & Observação registrada & Encaminhamento \\
\midrule
1º bim. & & & \\
2º bim. & & & \\
3º bim. & & & \\
4º bim. & & & \\
\bottomrule\end{tabular}\end{center}

\section{Roteiro de Reunião com a Família}

1) Acolher e explicar o método (o que o aluno já conquistou). 2) Mostrar o Registro de Progresso e as folhas feitas. 3) Explicar o que é rastreio escolar: observação pedagógica, nunca diagnóstico. 4) Indicar o profissional habilitado somente se sinais persistentes aparecerem. 5) Combinar rotina diária de 10 minutos de leitura compartilhada em casa. 6) Registrar os combinados por escrito.
